\section{Data Logging and Storage}
\label{sec:vayu_logging}

The data logging subsystem in Vayu provides persistent storage of system data for debugging, analysis, and future extensions such as flight data recording. It is built on top of the virtual file system (VFS) layer and is designed to ensure reliable operation without impacting real-time performance.

\subsection{File-Based Logging Architecture}

Logging is organized into multiple categories, each mapped to a dedicated file on storage. The current implementation supports separate logging streams for:
\begin{itemize}
    \item System-level events
    \item Communication (Navlink) data
    \item General-purpose logging
\end{itemize}

Each logging stream operates independently, allowing selective recording of different types of data without interference.

\subsection{Preallocation and Deterministic Writes}

To ensure predictable performance, log files are preallocated to fixed sizes during system initialization. This avoids dynamic file growth and reduces fragmentation, which is critical for maintaining consistent write latency in embedded storage systems.

Before logging begins, each file is expanded to its configured size and initialized, ensuring that all subsequent writes occur within a known memory region.

\subsection{Circular Logging Mechanism}

Logging is implemented using a circular buffer approach at the file level. Each log file maintains a write pointer that advances with each write operation. When the end of the file is reached, the write pointer wraps around to the beginning, overwriting older data.

\[
\text{write\_pos} = (\text{write\_pos} + \text{len}) \bmod \text{file\_size}
\]

This approach ensures continuous logging without requiring explicit file management or risking storage exhaustion.

\subsection{Thread-Safe Operation}

To support concurrent access from multiple tasks, each logging stream is protected using mutexes. This guarantees safe and consistent writes even when multiple subsystems attempt to log data simultaneously.

Write operations follow a structured sequence:
\begin{enumerate}
    \item Acquire mutex
    \item Seek to current write position
    \item Write data to storage
    \item Synchronize file to ensure persistence
    \item Update write pointer
    \item Release mutex
\end{enumerate}

This design ensures data integrity while maintaining simplicity.

\subsection{Integration with System Components}

The logging subsystem is integrated with other parts of Vayu, including communication and telemetry. Log data can be generated by various modules and stored persistently, while selected logs may also be transmitted in real time through the telemetry system.

Currently, the logging system is used to store calibration data and system-level information. These stored values can be retrieved across system restarts, enabling persistent configuration.

\subsection{Future Extensions}

The logging infrastructure is designed to support future expansion. Planned extensions include:
\begin{itemize}
    \item Full flight data recording (sensor, control, and state data)
    \item Blackbox-style logging for post-flight analysis
    \item Configurable logging levels and formats
    \item Efficient binary encoding for high-frequency data
\end{itemize}

With the underlying VFS and SD card support already in place, these features can be integrated without significant changes to the system architecture.

\subsection{System Behavior}

The logging subsystem operates independently of time-critical control tasks. By using preallocation, circular buffering, and mutex protection, it ensures that storage operations remain predictable and do not introduce blocking behavior in high-priority tasks.
